\chapter{绪论}

\section{研究背景}

在当今知识经济与"大科学"（Big Science）时代，科学研究日益超越传统的国家地理边界与政治壁垒，全球科学（Global Science）已不再是国家科学体系的简单叠加，而是一种高度联结、跨地域运行的知识生产体系，正在深刻重塑世界学术版图与科研资源配置格局\cite{limei2025, marginson2022}。科学家的跨国流动因此成为全球知识扩散、技术创新与合作网络构建的重要动力\cite{liu2024}。在这一宏观背景下，国际学术流动（international academic mobility）已成为当代高等教育与科学研究体系的重要组成部分。科研人员通过跨国流动嵌入不同制度环境与学术场域，在知识、技术与资源的跨境交换中实现能力提升与网络拓展。科学与技术人力资本（Scientific and Technical Human Capital, STHC）理论指出，科研人员的职业发展是技术资本与社会资本的获取和相互作用的过程\cite{canibano2019}，技术资本包含认知能力和隐性知识等，社会资本是专业与个人互动关系网络\cite{canibano2019}。

在世界科学体系的演变中，中国科学力量的崛起尤为引人关注。作为世界上最大的国际学生来源国之一，同时也是重要的国际学生接收国，中国在全球学术流动体系中的地位不断提升\cite{yang2022}。国家与全球科学体系之间并非简单从属关系，而是在互动中形成复杂的嵌入结构，尤其是在中国这类快速崛起的科研大国中，这种互动在学术资本的积累过程中发挥了独特的作用\cite{marginson2021}。Marginson指出，尽管中国的科研系统在全球科学体系中占据着日益重要的位置，但其与传统的欧美科学系统相比，依然保持一定的独立性。这种"嵌入式独立"结构在政策驱动型国际流动中表现更为明显\cite{marginson2022}。

我国长期将国际人才流动视为提升研究生培养质量与增强国家创新能力的重要战略工具。中共中央、国务院印发的《教育强国建设规划纲要（2024－2035年）》明确提出，要"完善教育对外开放战略策略，建设具有全球影响力的重要教育中心"，"提升全球人才培养和集聚能力。加强对出国留学人员的教育引导和服务管理。改革国家公派出国留学体制机制。鼓励支持选拔优秀人才到国际知名高校、研究机构研修，扩大国际学术交流和教育科研合作。"习近平总书记2025年5月31日在《求是》发表的重要文章《加快建设教育强国》中也强调，"学习借鉴国际先进经验是建设教育强国的重要途径"，要"深入推动教育对外开放，统筹'引进来'和'走出去'，不断提升教育国际影响力、竞争力和话语权"，这为科研人员出国留学、开展国际学术交流，以及研究生教育高质量发展提供了根本遵循。

自20世纪70年代末改革开放以来，中国政府高度重视通过公派出国留学推动科技人才培养与国际化战略。国家留学基金委员会（China Scholarship Council，CSC）自1996年成立以来\cite{csc_about}不断扩大资助规模，成为国家人才国际化战略的重要实施机构。其中，"国家建设高水平大学公派研究生项目"已成为规模最大的公派项目之一，每年资助大量博士研究生与博士后赴境外开展研究\cite{csc_2026guide}。CSC资助的博士生赴境外合作院校或科研机构开展研究，兼具国内学术训练与国际科研经历的双重属性，构成了本研究关注的核心群体。有研究指出，中国的科研人员，尤其是那些通过国家留学基金委（CSC）资助留学的博士生，往往能够在国际舞台上获得更高的学术声誉和更广泛的学术联系\cite{jianghua2023}。CSC资助的留学人员在海外学术环境中接受训练、开展合作研究，其学术资本积累过程与职业发展路径，成为观察国际学术流动影响的典型样本。

从职业阶段视角来看，欧洲委员会（EC）提出的科研职业阶段框架将研究人员划分为R1至R4四个阶段，其中博士研究生被明确界定为R1阶段（First Stage Researcher），即科研职业的起始期。这一阶段的研究者以系统科研训练为核心任务，在导师指导下逐步掌握研究方法、积累基础科研能力，并为未来的学术发展奠定基础。正是在这一关键起点阶段，国际流动经历对博士生的学术成长与职业发展产生的影响深远，值得深入探究。

\section{研究问题与研究意义}

\subsection{研究问题}

基于上述背景，国际学术流动对科研人员的影响已受到学界持续关注，既有研究在宏观层面积累了丰富的量化证据，证明流动经历与更高的论文产出、更广泛的国际合作之间存在显著关联。然而，这些研究呈现的更多是国际流动的可量化的结果指标，而对流动过程本身的内在机制关注相对有限。国际流动虽然被普遍认为能够扩大学术网络、提升科研能力并促进职业发展，但在不同制度环境、学科背景与个体发展阶段下，其具体作用机制仍较为复杂。科研人员在海外留学期间具体获取了哪些资源？这些资源又如何在回国后的科研活动与职业发展中发挥作用？对于这些问题，现有文献仍缺乏系统而深入的探讨。

从政策角度来看，中国国家留学基金委（CSC）每年投入巨额公共资源支持大规模的学术流动，然而目前的效用评价大多关注显性的宏观指标，如回国后的论文数量、职称晋升等，对于流动过程中科研人员所获得的隐性积累关注相对较少，因此，若要更全面地评估国际流动的真实效能，有必要从微观层面展开分析，尤其需要从个体视角探讨流动经历对科研人员学术资本积累的具体影响。

从个人角度来看，对于处于职业早期阶段的科研人员而言，海外经历并不必然等同于职业成功，一些海归博士在回国后可能面临"跨国资本溢价"与"本土网络脱节"的双重困境。博士阶段是科研人员职业发展的初始阶段，具有高度的不确定性，在这一阶段，国际学术流动带来的学术资本积累对个人职业路径的影响如何？如何能最大化程度地利用国际学术流动机会助力他们的职业发展？探讨这一问题，能够帮助我们从人才战略的角度理解如何更好地支持青年学者在职业初期建立竞争优势。

因此，考察CSC资助的博士生在国际流动过程中的学术资本积累，以及这些资本如何作用于职业发展，既是回应现实政策需要的必要之举，也是弥补既有理论研究缺口的学术要求。围绕"国际流动如何影响科研人员早期学术资本积累与职业发展"这一核心议题，本研究以CSC资助的博士研究生为具体观察对象，聚焦于流动经历的内在过程与微观机制，具体而言，本研究尝试回答以下三个递进式研究问题：

（1）国际学术流动过程中，科研人员主要积累了哪些类型的学术资本？

（2）国际学术流动特征如何影响科研人员的学术资本积累？

（3）国际学术流动所积累的学术资本对科研人员职业发展有何影响？

\subsection{研究意义}

本研究以国家留学基金委（CSC）资助的博士研究生为研究对象，采用深度访谈与扎根理论相结合的质性研究方法，通过对研究参与者国际流动经历的考察，探究科研人员在留学期间积累学术资本的类型、路径与差异，并追踪这些资本在回国后对其早期职业发展阶段的影响。本研究的价值体现在理论与实践两个层面。

从理论意义而言，首先，本研究拓展了布迪厄资本理论在国际学术流动情境下的适用性。布迪厄的资本理论为理解学术场域中的权力结构与资源分配提供了经典框架，但其原初概念主要基于法国社会的本土经验，对于国际流动情境下资本的生成、转化与利用问题关注有限。本研究通过对CSC资助博士生流动经历的深度分析，有望在经典资本理论框架的基础上，识别国际学术流动情境下学术资本的新维度，从而丰富和深化相关理论在国际学术流动领域的解释力。

其次，本研究补充了国际学术流动研究的微观机制视角。已有研究在宏观层面积累了大量关于流动效应的量化证据，但对流动过程内部的微观机制——即科研人员如何在日常科研实践中积累资本、网络如何在互动中逐步形成、隐性知识如何在师徒关系与同伴交流中悄然转移——缺乏深入的质性探索。本研究运用扎根理论方法对访谈文本进行系统分析，有望从行动者视角还原学术资本积累的过程性逻辑，为这一领域提供微观机制层面的理论补充。

最后，本研究回应中国情境下CSC资助博士生国际学术流动研究的理论空白。以往国际学术流动研究多以欧美国家的完整学位移动或博士后流动为主要对象，其理论框架与研究结论未必能直接迁移至中国情境。CSC资助的博士生在国内与境外学术体制之间流动，面临不同学术文化与身份认同的多重张力，这一独特的结构性位置可能催生有别于既有研究发现的资本积累路径。本研究以中国情境为出发点，有望为国际学术流动研究贡献来自全球南方与非西方经验的理论视角。

在实践意义方面，首先，本研究有助于为中国及其他发展中国家的教育和科研政策提供实证依据。通过分析中国科研人员在国际学术流动中的学术资本积累路径及其对职业发展的影响，本研究可以为政府优化公派留学政策、加强国际学术交流与合作提供政策建议，从而提升公派留学资源的配置效率与人才培养质量。其次，对于科研人员个体而言，本研究能够提供关于如何通过国际流动积累和运用学术资本的实用建议，帮助科研人员更好地规划学术生涯。特别是对即将申请或正在参与CSC公派项目的博士研究生，如何在有限的流动时期内最大化学术资本的积累，并在回国后有效地将其转化为职业发展的优势，是具有直接实践价值的现实问题。本研究通过呈现不同流动特征下资本积累与转化的成功路径与障碍因素，能够为科研人员的职业规划提供来自同侪经验的具体参照。此外，本研究还能够为高校与研究机构在培养国际化人才、优化人才培养模式等方面提供参考。

% TODO: 此处插入图表（原文标注"页码"）

\section{研究思路}

\subsection{研究内容}

本研究旨在探讨国际学术流动对科研人员早期学术资本积累与职业发展的影响机制。本文以国家留学基金委资助的博士研究生为研究对象，在系统梳理相关文献和理论的基础上，结合深度访谈资料，分析博士生在国际流动过程中积累的文化资本、社会资本、象征资本和心理资本，考察留学环境特征、个人特征和流动时长等因素对学术资本积累的影响，并进一步探讨这些资本在初职获得、职称晋升、项目申请和机会拓展中的具体作用。在此基础上，本文构建"国际流动特征---学术资本积累---职业发展"的分析框架，以呈现国际学术流动影响科研人员职业早期发展的过程与机制。最后，本文对全文的研究内容、研究方法和研究结论进行总结，并提出研究启示、不足与未来展望。本文主要包括六个章节，各章节内容概述如下：

第一章 绪论。本章主要介绍研究背景、研究问题与研究意义，说明研究思路、研究方法、研究贡献与论文结构，为全文奠定基础。

第二章 理论基础与文献综述。本章首先对国际学术流动与学术资本等核心概念进行界定，其次梳理布迪厄资本理论、社会资本理论和心理资本理论，并系统回顾国际学术流动与、科研产出及职业发展相关研究，在此基础上归纳既有研究的不足，明确本文的研究切入点。

第三章 研究设计。本章主要介绍访谈提纲设计、研究对象与资料收集过程，以及开放式编码、主轴编码和选择性编码等数据分析过程，说明本文如何运用扎根理论对访谈资料进行系统分析。

第四章 国际流动对科研人员学术资本积累的影响。本章首先分析博士生在国际学术流动中积累的学术资本类型，将其概括为文化资本、社会资本、象征资本与心理资本。随后，本章从留学环境特征、个人特征与流动时长三个方面，分析国际流动特征对学术资本积累的影响机制。

第五章 国际流动积累的学术资本对科研人员职业发展的影响。本章主要从初职获得、职称晋升、项目申请和机会拓展四个方面，分析不同类型学术资本在科研人员职业发展中的具体作用方式。

第六章 结论与展望。本章对全文研究发现进行总结，概括研究结论，提炼理论启示与实践启示，并指出研究不足与未来研究方向。

% TODO: 此处插入研究框架图
% \begin{figure}[htbp]
%   \centering
%   \includegraphics[width=\textwidth]{figures/research-framework.png}
%   \caption{研究框架图}
%   \label{fig:research-framework}
% \end{figure}

\subsection{研究方法}

本研究采用质性研究方法开展相关探究，原因主要有两点。第一，国际学术流动对科研人员早期学术资本积累与职业发展的影响，植根于博士生个体在国际流动过程中的具体互动情境、主观认知及其决策过程，因而呈现出较强的动态性、情境性与个体差异性。相较于其他研究方法，质性研究更有助于深入理解博士生在国际流动期间的实际经历与意义建构，从微观层面把握科研人员个体的具体体验及其行为逻辑，进而归纳早期学术资本形成的具体样态，并分析不同类型国际学术流动在资本积累上的差异。第二，现有研究虽然已对国际流动、学术资本积累与科研人员职业发展之间的关系有所讨论，但针对国家留学基金委资助博士生这一特定群体，尤其是围绕其国际流动经历如何影响早期学术资本积累与职业发展的问题，相关研究仍较为有限。因此，本研究尝试通过对访谈资料的持续比较与归纳分析，形成具有一定解释力的理论认识。扎根理论分析方法不预先设定明确的理论假设或因果路径，而是强调立足原始资料，对研究现象进行概念化、范畴化和理论整合。这一方法有助于从访谈材料中提炼国际学术流动影响科研人员早期学术资本积累与职业发展的内在逻辑，因而与本研究的探索性取向较为契合。

因此，本研究将采用质性研究方法，首先运用文献调研法梳理国际流动、学术资本积累与职业发展相关文献，厘清关键概念。其次采用访谈法，获得国家留学基金委（CSC）公派博士生海外科研经历、学术资本积累、回国后职业发展等文本数据。再运用扎根理论方法对访谈文本进行系统性、持续性分析，探索国际流动对于科研人员早期学术资本积累及职业发展的影响。主要研究方法的具体介绍如下：

（1）文献调研法

文献调研法是一种科学研究方法，通过调查文献获取资料，以全面了解和掌握研究问题。文献调研包括检索、查阅、整理、分析和评述现有相关文献资料的过程，本研究将选取中国知网、万方数据、Web of Science、Elsevier等文献数据库，系统梳理国际流动与学术资本、国际流动与职业发展、CSC资助留学人员的相关研究，对现有文献进行深入的整合与评述，为后续研究构建扎实的理论框架。

（2）访谈法

访谈法是一种通过访员与受访者之间的问答互动来搜集数据的研究手段，在相对开放、交流性的对话情境中，受访者往往能够更充分地表达自身经验与观点，提供更为丰富的信息，还可能补充关键细节与相关背景，呈现出那些难以被量化指标所捕捉的个人化认知与深层体验。访谈法的优点是能够获取主观体验、行为动机及复杂社会背景下的深层逻辑，适用于探索性研究。访谈法因易于操作、成本可控且重视受访者的主体体验，在图书情报领域的质性研究中，是应用最广泛的资料收集方法\cite{baoxin2021}。半结构化访谈（semi-structured interview）是质性研究常见的一种访谈方式，又称深度访谈（in-depth interview）\cite{sunxiaoe2012}，它使用预先设计的访谈提纲作为框架，在执行过程中允许研究者根据受访者的回答灵活调整问题顺序、深入探索新主题或补充追问，给予受访者和研究者一定的自由度，使数据收集既有系统性又有深度。因此，本研究将采用半结构化访谈方式，一对一与受访者进行深度访谈，在访谈过程中基于预设框架开展灵活追问，给予受访者充分的空间来思考和回答，挖掘受访者的真实体验、感知与反思，采集一手叙事资料，为后续扎根理论分析提供数据支撑。本研究关注的对象不是正处于国际学术流动阶段的博士生，而是已完成访学并任职于高校或其他科研单位的博士群体，因此一方面，从个人层面而言，受访者在经历多年职业实践后对过往访学经历进行回溯性反思，往往能产生更具深度与价值的认知，另一方面，基于长期记忆的反思，不可避免地存在记忆模糊、细节缺失或偏差等风险，这也构成了本研究的一项局限性。

（3）扎根理论

扎根理论（Grounded Theory）是一种质性研究方法，由巴尼·格拉泽（Barney Glaser）和安塞尔姆·施特劳斯（Anselm Strauss）于1967年首次提出，其特点是以文字叙述为材料、以归纳法为论证步骤、以构建主义为前提，强调自下而上地、从经验资料中生成理论\cite{chenxiangming1999}。扎根理论方法要求研究者对理论保持敏感，在适当使用已有文献的基础上，研究者要将资料与个人解释结合，构建起资料、个人解释与前人成果间的三角互动关系，最终形成理论。扎根理论适用于理论体系尚不完善的领域，因此契合本研究探索需求。

扎根理论分为三个流派：Glaser和Strauss提出的经典扎根理论、Strauss和Corbin提出的程序化扎根理论以及Charmaz提出的建构扎根理论。经典流派强调理论从数据中自然涌现，通过持续比较分析避免预设框架，侧重数据驱动的自然生成性，避免研究者的干预；程序化流派采用系统化三级编码（开放式、主轴式、选择性）确保规范性和可重复性，强调高度结构化的操作流程，追求方法论严谨性；建构型流派主张理论是研究者与数据在社会背景下的共同建构，注重主观解释和意义协商，突出研究者的主动参与和情境化分析。其中，程序化扎根理论在中国的应用最为广泛，也是信息资源管理领域普遍采取的数据分析方式\cite{zhangxiangyihong2025}，因此，本研究将主要依据Strauss和Corbin所提出的程序化扎根理论方法，运用其系统的三级编码程序（开放式、主轴式、选择性编码）对资料进行分析。

\section{研究贡献}

本研究的贡献主要体现在以下三个方面：

第一，在理论框架上，本文构建了涵盖文化资本、社会资本、象征资本与心理资本四个维度的学术资本理论框架，并将其置于"国际流动特征---学术资本积累---职业发展"的分析链条之中，丰富了现有国际流动理论。相较于既有研究较多关注论文产出、国际合作或职业晋升等结果性指标，本文进一步关注国际学术流动影响职业发展的中间过程，有助于从过程视角理解国际学术流动对科研人员早期发展的作用机制。

第二，在研究内容上，本文明晰了国际流动特征对学术资本积累的影响机制。基于访谈资料，本文从留学环境特征、个人特征与流动时长三个方面，分析了国际流动特征对学术资本积累的影响，揭示了不同流动条件下学术资本积累的差异与作用路径。

第三，在研究发现上，本文探究了国际流动积累的学术资本对科研人员早期职业发展的作用机制。通过考察不同类型学术资本在初职获得、职称晋升、项目申请和机会拓展中的具体作用，揭示了学术资本转化为职业发展优势的差异化路径。

第四，在研究对象与经验情境上，本文聚焦于成长型科研人员——国家留学基金委资助的博士研究生。对这一群体的考察，有助于丰富中国情境下国际学术流动研究的经验材料，并为理解CSC资助博士生的国际流动经历及其影响提供具体的实证依据。
